% !TEX program = xelatex
\documentclass[11pt,a4paper]{ctexart}

\usepackage{amsmath,amssymb,amsfonts}
\usepackage{booktabs}
\usepackage{geometry}
\usepackage{hyperref}
\usepackage{enumitem}
\usepackage{xcolor}
\usepackage{longtable}
\usepackage{array}

\geometry{left=2.5cm,right=2.5cm,top=2.5cm,bottom=2.5cm}
\hypersetup{colorlinks=true,linkcolor=blue,urlcolor=blue}

\newcolumntype{L}[1]{>{\raggedright\arraybackslash}p{#1}}

\title{%
  \textbf{星闪 SLB 物理层}\\[0.4em]
  \Large 6.2.2 基础帧结构和物理资源技术文档\\[0.3em]
  \large CP-OFDM 波形 \textbullet\ 超帧/无线帧 \textbullet\ COT/TTI \textbullet\ 信道载波 \textbullet\ RE 网格
}
\author{依据 T/XS 10001-2025《星闪无线通信系统 接入层\\
同步低时延宽带空口 SLB 技术要求和测试方法》整理\\
（团体标准 20250326 版）}
\date{2026 年 7 月 6 日}

\begin{document}
\maketitle
\tableofcontents
\newpage

%======================================================================
\part{总览}
%======================================================================

\section{文档范围与模块划分}

本文档对应 SLB 120\,kHz 物理层协议第 6.2.2 节``基础帧结构和物理资源''，涵盖 6.2.2.1\textasciitilde 6.2.2.10 全部子节。该节在 6.2.1 内容分类之上，定义 120\,kHz CP-OFDM 系统的\textbf{时频资源坐标系}——波形参数、帧层次、COT/TTI、载波结构、RE 索引及通信域/直通链路资源边界，是 6.2.3\textasciitilde 6.2.5 进行信号/控制/数据映射的前置基础。

\begin{table}[h]
\centering
\small
\begin{tabular}{L{2.2cm}L{4.8cm}L{6.5cm}}
\toprule
\textbf{子节号} & \textbf{子节主题} & \textbf{核心输出/行为} \\
\midrule
6.2.2.1 & CP-OFDM 波形与五种符号类型 & $\Delta f$、CP/GAP 长度、边界符号规则（表 1） \\
6.2.2.2 & 超帧/无线帧与链路符号 & 1\,ms 超帧、$N_{\mathrm{sym}}^{\mathrm{frame}}$、G/T/A/D 符号配置（表 2/3） \\
6.2.2.3 & COT 与 TTI & 占信道时长、$2^{N_{\mathrm{TTI}}}$ 无线帧块、链路时分规则 \\
6.2.2.4 & 信道与载波 & 20\,MHz 基础信道、161 子载波、$N_{\mathrm{CH}}$/$L_{\mathrm{CH}}$（表 4） \\
6.2.2.5 & 天线端口 & 逻辑端口定义；SAB/G-PCIB/广播同端口 \\
6.2.2.6 & 资源网格 & $161 \times N_{\mathrm{sym}}^{\mathrm{frame}}$ 时频平面 \\
6.2.2.7 & 资源元素 RE & $(k,l)_p$ 索引、$a_{k,l}$ 映射、DC/空闲约束 \\
6.2.2.8 & 通信域 & G 节点工作载波上的可调度资源集合 \\
6.2.2.9 & G/T/附链路 & 通信域内三条逻辑链路的方向与符号类型 \\
6.2.2.10 & 直通链路 & 竞争获取的 D 符号资源；先发/后发节点 \\
\bottomrule
\end{tabular}
\end{table}

\section{与相关章节的关系}

\begin{table}[h]
\centering
\small
\begin{tabular}{L{2.2cm}L{2.2cm}L{9.1cm}}
\toprule
\textbf{相关节号} & \textbf{关系} & \textbf{说明} \\
\midrule
6.1 & 上游 & $T_s$、$f_s$、比特序约定；CP-OFDM 波形总体框架 \\
6.2.1 & 上游 & 信号/控制/数据/PIB 分类；6.2.2 提供映射坐标 \\
6.2.3 & 下游 & FTS/STS、DMRS 等在 RE $(k,l)$ 上的映射 \\
6.2.4/6.2.5 & 下游 & SAB/RAB/G-PCIB 块内符号 \#3/\#5 等索引依赖帧结构 \\
6.2.9/6.2.10 & 下游 & OFDM 合成与射频发射使用符号时长与 CP 参数 \\
6.5.3 & 下游 & 符号号 $l$ 与循环前缀类型 $N_{\mathrm{CP}}$ 依赖表 1/表 3 \\
8.3 & 上游 & 基础信道中心频率与多载波频域布局 \\
6.3 & 平行 & 灵活带宽在相同帧结构上的 $\alpha$ 缩放 \\
\bottomrule
\end{tabular}
\end{table}

\section{基本时间与频率量}

\begin{table}[h]
\centering
\small
\begin{tabular}{llll}
\toprule
\textbf{基本量} & \textbf{定义/公式} & \textbf{数值} & \textbf{标准条款} \\
\midrule
基准频率 $f_s$ & 物理层采样/定时基准 & 30.72\,MHz & 6.1 \\
基本时间单位 $T_s$ & $T_s = 1/f_s$ & $\approx 32.55$\,ns & 6.1 \\
子载波间隔 $\Delta f$ & $\Delta f = f_s/256$ & 120\,kHz & 6.2.2.1 \\
OFDM 有效部分 & 不含 CP 的数据段 & $256 T_s \approx 8.333$\,$\mu$s & 6.2.2.1 \\
超帧时长 $T_f$ & $T_f = 30720 \times T_s$ & 1\,ms & 6.2.2.2 \\
\bottomrule
\end{tabular}
\end{table}

\newpage
%======================================================================
\part{6.2.2 基础帧结构和物理资源（工程解读）}
%======================================================================

\section{协议章节对应}

本节对应 T/XS 10001-2025 第 6.2.2 节，完整涵盖 6.2.2.1\textasciitilde 6.2.2.10 共 10 个子节。

\section{功能定义}

6.2.2 的核心功能为：定义 120\,kHz CP-OFDM 系统的\textbf{时频资源框架}，使发送/接收端能够：
\begin{enumerate}[nosep]
  \item \textbf{确定符号物理参数}：按符号类型查表 1 获取 CP/GAP 长度与符号总时长；
  \item \textbf{定位时间资源}：在超帧/无线帧/符号三级时间轴上索引 PIB 与业务资源；
  \item \textbf{索引频域资源}：在节点工作载波上按 $k$、子载波组、工作子载波组编号；
  \item \textbf{映射 RE}：在 $161 \times N_{\mathrm{sym}}^{\mathrm{frame}}$ 资源网格上以 $(k,l)_p$ 写入 $a_{k,l}$；
  \item \textbf{区分域与链路}：通信域（G/T/A 符号）与直通链路（D 符号）及 COT/TTI 占用规则。
\end{enumerate}

\section{模块边界（上下游及职责划分）}

\subsection{上游模块}

\begin{table}[h]
\centering
\small
\begin{tabular}{L{4cm}L{4.5cm}L{5.5cm}}
\toprule
\textbf{上游} & \textbf{输入} & \textbf{说明} \\
\midrule
6.1 物理层一般描述 & $T_s$、$f_s$、比特顺序约定 & 全局 PHY 基准 \\
高层/MAC 配置 & 符号类型、$N_{\mathrm{TTI}}$、$N_{\mathrm{CH}}$、$L_{\mathrm{CH}}$ & 帧长与带宽配置 \\
8.3 信道号与中心频率 & 基础信道中心频率 & 多载波频域布局 \\
\bottomrule
\end{tabular}
\end{table}

\subsection{下游模块}

\begin{table}[h]
\centering
\small
\begin{tabular}{L{4cm}L{4.5cm}L{5.5cm}}
\toprule
\textbf{下游} & \textbf{输出} & \textbf{说明} \\
\midrule
6.2.3 物理层信号 & RE $(k,l)$ 时频位置 & FTS/STS、DMRS、CSI-RS 等映射基准 \\
6.2.4/6.2.5 控制/数据 & 资源网格与符号 \# 索引 & PIB 块内 RE 分配 \\
6.2.9 基带信号生成 & 符号时长与 CP 参数 & OFDM 调制链输入 \\
6.5.3 QPSK 序列 & $n$、$l$、$N_{\mathrm{ID}}$、$N_{\mathrm{CP}}$ & 符号号 $l$ 与 CP 类型对应 \\
\bottomrule
\end{tabular}
\end{table}

\subsection{本模块不负责的内容}

波形调制算法（6.2.9/6.2.10）、具体 PIB 比特字段（6.2.4）、信号序列生成（6.2.3）、信道编码（6.2.6）、射频指标（第 8 章）及 MAC 调度策略。

\section{在 SLB 通信流程中的角色}

6.2.2 处于 PHY 层\textbf{资源坐标定义层}：在 6.2.1 内容分类之后、6.2.3\textasciitilde 6.2.5 具体映射之前。

\textbf{接收端资源定位故事线：}
\begin{enumerate}[nosep]
  \item \textbf{帧同步}：锁定超帧/无线帧边界（依赖 6.2.3 FTS/STS，坐标系来自本节）；
  \item \textbf{符号定时}：按符号类型与表 3 区分一般/边界 CP-OFDM 符号；
  \item \textbf{载波选择}：确定节点工作信道/载波及基础载波序号；
  \item \textbf{RE 解调}：在 $(k,l)_p$ 上读取 $a_{k,l}$，区分 G/T/A/D 符号上下文；
  \item \textbf{块解析}：结合 6.2.1 PIB 分类，在已定位 RE 上解析控制/数据。
\end{enumerate}

\textbf{坐标轴关系：}
\begin{itemize}[nosep]
  \item \textbf{时间轴}：超帧 (1\,ms) $\rightarrow$ 8 个无线帧 $\rightarrow$ $N_{\mathrm{sym}}^{\mathrm{frame}}$ 个符号 $\rightarrow$ RE 时域索引 $l$；
  \item \textbf{频域轴}：节点工作信道 ($N_{\mathrm{CH}}$ 个 20\,MHz) $\rightarrow$ 节点工作载波 $\rightarrow$ 161 子载波/基础载波 $\rightarrow$ RE 频域索引 $k$。
\end{itemize}

%----------------------------------------------------------------------
\section{关键参数与强制要求（提炼版）}
%----------------------------------------------------------------------

\subsection{CP-OFDM 波形与符号类型（6.2.2.1，表 1）}

物理层信号、控制、数据均基于 $\Delta f = 120$\,kHz CP-OFDM 传输。有效数据段长度恒为 $256 T_s$；不含 CP 的 OFDM 符号满足 $T_{\mathrm{CP}}=0$，$T_{\mathrm{Symb}}^{\mathrm{OFDM}} = 256 T_s$。

\begin{table}[h]
\centering
\caption{表 1\quad 120\,kHz CP-OFDM 符号设计（提炼）}
\small
\begin{tabular}{L{1.8cm}lllll}
\toprule
\textbf{符号类型} & \textbf{一般 CP} & \textbf{一般符号长} &
\textbf{GAP1} & \textbf{边界符号长} & \textbf{边界符号} \\
\midrule
类型一 A & $18T_s$ & $274T_s$ & $20T_s$ & $276T_s$ & 有（\#0,\#7） \\
类型一 B & $18T_s$ & $274T_s$ & $22T_s$ & $278T_s$ & 有（\#0） \\
类型二 & $39T_s$ & $295T_s$ & $44T_s$ & $300T_s$ & 有（\#0） \\
类型三 & $64T_s$ & $320T_s$ & --- & --- & 无 \\
类型四 & $128T_s$ & $384T_s$ & --- & --- & 无 \\
\bottomrule
\end{tabular}
\end{table}

\textbf{与 6.5.3 衔接}：通信域 CP-OFDM 符号类型 $\rightarrow$ $N_{\mathrm{CP}}$：类型一 A/B $\rightarrow$ 3；类型二 $\rightarrow$ 2；类型三 $\rightarrow$ 1；类型四 $\rightarrow$ 0。

\subsection{超帧/无线帧与符号配置（6.2.2.2，表 2/3）}

\begin{table}[h]
\centering
\small
\begin{tabular}{lll}
\toprule
\textbf{参数} & \textbf{取值/规则} & \textbf{标准条款} \\
\midrule
超帧时长 $T_f$ & $30720 T_s = 1$\,ms & 6.2.2.2 \\
超帧编号 & \#0\textasciitilde\#65535 循环 & 6.2.2.2 \\
每超帧无线帧数 & 8（\#0\textasciitilde\#7） & 6.2.2.2 \\
$N_{\mathrm{sym}}^{\mathrm{frame}}$ & 类型一 A/B: 14；二: 13；三: 12；四: 10 & 表 2 \\
通信域符号 & G/T/A 符号 + 切换间隔 & 6.2.2.2 \\
直通链路符号 & D 符号 + 切换间隔 & 6.2.2.2 \\
\bottomrule
\end{tabular}
\end{table}

\begin{table}[h]
\centering
\caption{表 3\quad 无线帧符号配置（边界 CP 位置）}
\small
\begin{tabular}{llll}
\toprule
\textbf{符号类型} & \textbf{符号数} & \textbf{一般 CP 符号 $l$} & \textbf{边界 CP 符号 $l$} \\
\midrule
类型一 A & 14 & \#1\textasciitilde\#6, \#8\textasciitilde\#13 & \#0, \#7 \\
类型一 B & 14 & \#1\textasciitilde\#13 & \#0 \\
类型二 & 13 & \#1\textasciitilde\#12 & \#0 \\
类型三 & 12 & \#0\textasciitilde\#11 & --- \\
类型四 & 10 & \#0\textasciitilde\#9 & --- \\
\bottomrule
\end{tabular}
\end{table}

\subsection{COT 与 TTI（6.2.2.3）}

\begin{table}[h]
\centering
\small
\begin{tabular}{lll}
\toprule
\textbf{参数/规则} & \textbf{取值/约束} & \textbf{标准条款} \\
\midrule
$N_{\mathrm{TTI}}$ & 0\textasciitilde 6 & 6.2.2.3 \\
TTI 无线帧数 & $2^{N_{\mathrm{TTI}}}$（1/2/4/8/16/32/64） & 6.2.2.3 \\
COT 定义 & 竞争成功后首帧至释放前末帧，整数个无线帧 & 6.2.2.3 \\
边界对齐 & COT/TTI 边界与超帧边界\textbf{无固定关系} & 6.2.2.3 \\
通信域 TTI & 至少含 G 链路且 G 在 TTI 开头 & 6.2.2.3 \\
链路切换 & 不同类型链路间插入 1 符号切换间隔 & 6.2.2.3 \\
TTI 末符号 & 必须为切换间隔 & 6.2.2.3 \\
\bottomrule
\end{tabular}
\end{table}

\subsection{信道、载波与子载波组（6.2.2.4，表 4）}

\begin{table}[h]
\centering
\small
\begin{tabular}{lll}
\toprule
\textbf{参数} & \textbf{取值/规则} & \textbf{标准条款} \\
\midrule
基础信道带宽 & 20\,MHz & 6.2.2.4 \\
每基础载波子载波数 & 161；\#80 为 DC，不映射 & 6.2.2.4 \\
有效子载波 & 除 \#80 外均可映射 & 6.2.2.4 \\
基础子载波组 & 每 10 个有效子载波 = 1 组，共 16 组 & 6.2.2.4 \\
分隔子载波 & 频域相邻基础载波间，不映射 & 6.2.2.4 \\
节点工作信道 & 连续 $N_{\mathrm{CH}}$ 个基础信道并集 & 6.2.2.4 \\
工作子载波组编号 & \#0\textasciitilde\#$(\lfloor 16 N_{\mathrm{CH}}/L_{\mathrm{CH}} \rfloor - 1)$ & 6.2.2.4 \\
\bottomrule
\end{tabular}
\end{table}

\begin{table}[h]
\centering
\caption{表 4\quad $N_{\mathrm{CH}}$ 与 $L_{\mathrm{CH}}$ 对应关系}
\small
\begin{tabular}{c*{18}{c}}
\toprule
$N_{\mathrm{CH}}$ &
1 & 2 & 3 & 4 & 5 & 6 & 7 & 8 & 9 & 10 & 11 & 12 & 13 & 14 & 15 & 16 & 25 & 32 \\
\midrule
$L_{\mathrm{CH}}$ &
1 & 2 & 4 & 4 & 8 & 8 & 8 & 8 & 16 & 16 & 16 & 16 & 16 & 16 & 16 & 16 & 32 & 32 \\
\bottomrule
\end{tabular}
\end{table}

%----------------------------------------------------------------------
\section{本节核心详解}
%----------------------------------------------------------------------

\subsection{6.2.2.1 CP-OFDM 波形}

\textbf{功能}：为全部物理层内容提供统一时域符号结构。

\textbf{输入输出}：输入符号类型配置 $\rightarrow$ 输出每符号 CP 长度、总时长、是否边界符号。

\textbf{上下文}：通信域/直通链路均使用 120\,kHz SCS；类型三/四 CP 更长，适用于特定时延/覆盖场景；边界符号用于无线帧首尾（表 3），影响 6.5.3 中 $N_{\mathrm{CP}}$ 查表。

\subsection{6.2.2.2 超帧与无线帧}

\textbf{功能}：建立周期性时间网格，承载 G/T/A/D 链路符号序列。

\textbf{输入输出}：输入符号类型 $\rightarrow$ 输出 $N_{\mathrm{sym}}^{\mathrm{frame}}$、符号 \#0\textasciitilde\#$N{-}1$ 及 CP 类型。

\textbf{上下文}：6.2.4 中 SAB \#0\textasciitilde\#4、广播 \#5/\#7 等符号索引均引用本节无线帧符号编号；类型一 A 的 \#7 为边界符号，与 SAB/G-PCIB 典型布局相关。

\subsection{6.2.2.3 COT 与 TTI}

\textbf{功能}：描述一次信道占用内的传输时间块，支撑非连续/连续传输模式调度。

\textbf{输入输出}：输入 $N_{\mathrm{TTI}}$、传输模式 $\rightarrow$ 输出 TTI 无线帧数、COT 内 TTI 顺序号。

\textbf{上下文}：竞争接入后 COT 起始无线帧不必对齐超帧；通信域 TTI 必须以 G 链路开头，保证 SAB/G-PCIB 先于 T/A 链路；附链路/直通链路多发送节点间亦需 1 符号切换间隔。

\subsection{6.2.2.4 信道与载波}

\textbf{功能}：定义频域资源层次，支撑多载波聚合与调度资源指示。

\textbf{频域结构示意}：
\begin{itemize}[nosep]
  \item 单基础载波：\#0\textasciitilde\#79 $|$ \#80(DC) $|$ \#81\textasciitilde\#160；
  \item 16 个基础子载波组/载波；$N_{\mathrm{CH}}$ 个载波聚合为节点工作载波；
  \item 多载波间分隔子载波不承载内容（见标准图 51）。
\end{itemize}

\textbf{上下文}：6.2.4 DSCI 中``基础载波序号''、``工作子载波组''等字段均基于本节编号；$N_{\mathrm{CH}}$、$L_{\mathrm{CH}}$ 必须按表 4 配对。

\subsection{6.2.2.5\textasciitilde 6.2.2.7 天线端口、资源网格与 RE}

\textbf{天线端口}（6.2.2.5）：逻辑概念——同端口相邻符号信道可互推。通信域 SAB、G-PCIB、广播、G-DMRS-C 同端口；直通链路同节点发送内容同端口。

\textbf{资源网格}（6.2.2.6）：每天线端口 $p$、每基础载波、每无线帧 $\rightarrow$ $161 \times N_{\mathrm{sym}}^{\mathrm{frame}}$ RE 矩阵。

\textbf{RE}（6.2.2.7）：最小映射单元 $(k,l)_p$，$k{=}0,\ldots,160$，$l{=}0,\ldots,N_{\mathrm{sym}}^{\mathrm{frame}}{-}1$，映射 $a_{k,l}^{(p)}$；空闲 RE 映射 0；$a_{80,l}\equiv 0$ 恒成立。

\subsection{6.2.2.8\textasciitilde 6.2.2.10 通信域、三链路与直通链路}

\begin{table}[h]
\centering
\small
\caption{通信域与直通链路对比}
\begin{tabular}{L{2.5cm}L{4.5cm}L{3.5cm}L{3cm}}
\toprule
\textbf{概念} & \textbf{资源范围/获取方式} & \textbf{符号类型} & \textbf{标准条款} \\
\midrule
通信域 & G 节点工作载波：SAB、G-PCIB、广播、可调度资源 & G/T/A & 6.2.2.8 \\
G 链路 & G$\rightarrow$至少一个 T & G 符号 & 6.2.2.9 \\
T 链路 & T$\rightarrow$G & T 符号 & 6.2.2.9 \\
附链路 & T$\rightarrow$T 和/或 G & A 符号 & 6.2.2.9 \\
直通链路 & 竞争信道；先发/后发节点 & D 符号 & 6.2.2.10 \\
\bottomrule
\end{tabular}
\end{table}

%----------------------------------------------------------------------
\section{数据类型、长度与维度}
%----------------------------------------------------------------------

\begin{table}[h]
\centering
\small
\begin{tabular}{llll}
\toprule
\textbf{对象} & \textbf{维度/规模} & \textbf{单位} & \textbf{标准条款} \\
\midrule
子载波索引 $k$ & 0\textasciitilde 160（\#80 恒零） & 子载波 & 6.2.2.7 \\
符号索引 $l$ & 0\textasciitilde $N_{\mathrm{sym}}^{\mathrm{frame}}{-}1$ & 符号 & 6.2.2.2 \\
$N_{\mathrm{sym}}^{\mathrm{frame}}$ & 10/12/13/14 & 符号/无线帧 & 表 2 \\
超帧内无线帧 & 8 & 无线帧/超帧 & 6.2.2.2 \\
资源网格 RE 数 & $161 \times N_{\mathrm{sym}}^{\mathrm{frame}}$ & RE/帧/载波/端口 & 6.2.2.6 \\
基础子载波组/载波 & 16 & 组 & 6.2.2.4 \\
$N_{\mathrm{CH}}$ & 1,2,3,4,5,6,7,8,9,10,11,12,13,14,15,16,25,32 & 基础信道 & 表 4 \\
TTI 无线帧数 & $2^{N_{\mathrm{TTI}}}$，$N_{\mathrm{TTI}}\in\{0,\ldots,6\}$ & 无线帧 & 6.2.2.3 \\
$a_{k,l}^{(p)}$ & 复数值 & --- & 6.2.2.7 \\
\bottomrule
\end{tabular}
\end{table}

%----------------------------------------------------------------------
\section{时序关系与触发条件}
%----------------------------------------------------------------------

\begin{enumerate}[nosep]
  \item \textbf{超帧周期}：$T_f = 1$\,ms，编号 \#0\textasciitilde\#65535 循环，与业务调度周期独立；
  \item \textbf{无线帧}：超帧内固定 8 帧，符号类型决定每帧符号数；
  \item \textbf{COT 触发}：非连续模式下竞争信道成功后从首无线帧起算；
  \item \textbf{TTI 划分}：COT 或连续模式下每 $2^{N_{\mathrm{TTI}}}$ 无线帧为一个 TTI；
  \item \textbf{链路时分}：通信域 TTI 内 G 链路在首段，G/T/A 间以 1 符号切换间隔分隔；
  \item \textbf{TTI 结束}：末符号固定为切换间隔，下一 TTI 或释放信道；
  \item \textbf{直通链路}：竞争获取 COT 后配置 D 符号序列，无 G 节点调度前置。
\end{enumerate}

%----------------------------------------------------------------------
\section{处理步骤（接收侧资源定位）}
%----------------------------------------------------------------------

\textbf{Step 1}：读取广播/同步信息中的符号类型，查表 2 得 $N_{\mathrm{sym}}^{\mathrm{frame}}$。

\textbf{Step 2}：完成帧同步后，按超帧边界划分 8 个无线帧，每帧内按表 3 标记边界/一般 CP 符号。

\textbf{Step 3}：根据 $N_{\mathrm{TTI}}$ 与 COT 起始帧，计算当前 TTI 顺序号及所含无线帧范围。

\textbf{Step 4}：解析无线帧符号配置，区分 G/T/A/D 符号段与切换间隔符号。

\textbf{Step 5}：确定节点工作载波（$N_{\mathrm{CH}}$）及目标基础载波序号。

\textbf{Step 6}：在资源网格上按 $(k,l)_p$ 索引 RE，校验 $k{\neq}80$ 且 $a_{k,l}{\neq}0$ 后方进入 6.2.3/6.2.4 解调。

%----------------------------------------------------------------------
\section{约束与实现要点}
%----------------------------------------------------------------------

\begin{enumerate}[nosep]
  \item \textbf{DC 子载波}：$a_{80,l}\equiv 0$，任何信号/控制/数据均不得映射至 \#80；
  \item \textbf{分隔子载波}：多载波频域相邻时，分隔子载波不映射任何内容；
  \item \textbf{边界 CP 位置}：必须严格按表 3 配置，否则符号定时与 6.5.3 的 $l$/$N_{\mathrm{CP}}$ 不一致；
  \item \textbf{TTI 结构}：通信域 TTI 末符号必须为切换间隔；G 链路必须在 TTI 开头；
  \item \textbf{边界不对齐}：COT/TTI 起始可与超帧任意对齐，实现须支持任意无线帧边界；
  \item \textbf{表 4 配对}：$N_{\mathrm{CH}}$ 与 $L_{\mathrm{CH}}$ 必须按表 4 取值，否则工作子载波组编号错误；
  \item \textbf{空闲 RE}：未分配 RE 映射 $a_{k,l}{=}0$，接收端不得当作有效载荷；
  \item \textbf{天线端口}：通信域 SAB/G-PCIB/广播/G-DMRS-C 须使用同一天线端口做信道估计；
  \item \textbf{域不可混用}：通信域 G/T/A 符号配置不适用于直通链路 D 符号场景；
  \item \textbf{OFDM 简称}：无 CP 的 OFDM 符号 $T_{\mathrm{CP}}{=}0$；与 CP-OFDM 统称``符号''时须区分上下文。
\end{enumerate}

%----------------------------------------------------------------------
\section{与标准其他章节的衔接}
%----------------------------------------------------------------------

\begin{itemize}[nosep]
  \item \textbf{6.2.1}：PIB 与内容分类；6.2.2 提供映射坐标；
  \item \textbf{6.2.3}：FTS/STS/DMRS 等在 RE $(k,l)$ 上的映射规则；
  \item \textbf{6.2.4/6.2.5}：控制/数据在 G/T/A/D 符号上的块内布局与符号 \# 索引；
  \item \textbf{6.5.3}：符号号 $l$ 与 $N_{\mathrm{CP}}$ 依赖表 1/表 3 符号类型；
  \item \textbf{6.2.9/6.2.10}：基带 OFDM 合成与射频发射使用本节符号时长参数；
  \item \textbf{8.3}：基础信道中心频率与多载波频域布局；
  \item \textbf{6.3}：灵活带宽在相同帧结构上的带宽扩展。
\end{itemize}

\vfill
\begin{center}
\small\textcolor{gray}{字段级定义与完整参数以 T/XS 10001-2025 第 6 章及 6.2.2 各子节为准；本文档提供 120\,kHz 帧结构与 RE 坐标系工程解读，供 SLB 物理层实现与联调参考。}
\end{center}

\end{document}
